\section{The Partisan Distribution of Ensemble Plans}
\label{sec:bell-curve}

\subsection{The Bell Curve Around Geography}

The central insight of ensemble redistricting research is that, holding
geography fixed, the distribution of partisan seat outcomes across valid plans
is approximately bell-shaped.  The mean of this distribution is determined
by geography: the spatial concentration of partisan voters.  The width is
determined by the flexibility of the districting plan space: how many valid
plans exist, and how diverse their partisan outcomes are.

More formally, let $S_D(P)$ denote the number of Democratic-leaning districts
in plan $P$ under a fixed electoral baseline.  For a \recom{} ensemble
$\mathcal{E}$ on a state with $k$ seats:

\begin{equation}
  S_D \sim \text{approximately } \mathrm{Binomial}(k, p_{\text{geo}})
  \text{ with overdispersion},
  \label{eq:partisan-dist}
\end{equation}

where $p_{\text{geo}}$ is the geographically determined probability that
a random valid district leans Democratic.  For highly competitive states
(Wisconsin, Pennsylvania), $p_{\text{geo}} \approx 0.50$.  For sorted
states (Georgia), $p_{\text{geo}} \approx 0.36$.

The distribution is approximately symmetric around the geographic mean for
competitive states and skewed for sorted states.  Crucially, the distribution
is \emph{not} centered on proportionality: in sorted states, even a perfectly
random plan tends to produce more Republican seats than the statewide vote
share would imply.

\subsection{Width of the Distribution}

The standard deviation $\sigma_{S_D}$ of the partisan seat distribution
depends on:

\begin{itemize}
  \item $k$ (more districts $\to$ wider distribution in absolute terms, narrower
    in relative terms);
  \item Geographic compactness of the state (more irregular states $\to$
    wider distribution);
  \item Degree of partisan sorting (highly sorted states $\to$ narrower
    distribution near the geographic mean).
\end{itemize}

Empirically, $\sigma_{S_D} \approx 1.0$--$1.5$ seats for $k = 8$--$17$
and $\sigma_{S_D} \approx 2.0$--$3.5$ for $k = 38$--$52$.  A difference
of two seats from the ensemble median is thus within one to two standard
deviations for most states.

\subsection{What the Distribution Does Not Show}

The partisan distribution of ensemble plans does \emph{not} show:
\begin{itemize}
  \item The ``correct'' partisan outcome for the state;
  \item What proportionality requires;
  \item Whether any particular plan within the distribution is
    constitutionally required.
\end{itemize}

It shows what outcomes are geographically plausible.  A plan at the 99th
percentile is implausible without intent.  A plan at the 54th percentile
is typical.  Whether typical is fair is a separate normative question,
addressed in the context of proportionality in Section~\ref{sec:corridor}.
